\section*{1. System Model, Latent Uncertainty, and Decision Objective}

\textbf{Swing equation (IEEE-14, network-coupled, 2nd-order Kuramoto):}
\[
\dot{\theta}_i(t)=\omega_i(t)
\]
\[
M\,\dot{\omega}_i(t)=
P_{m,i}
-\sum_{j=1}^{N}B_{ij}\sin(\theta_i(t)-\theta_j(t))
-(D+K)\omega_i(t)
+u^{\text{probe}}_{\xi,i}(t)
+u^{\text{ctrl}}_{\gamma,i}(t)
\]

\textbf{Latent parameters (fixed per episode):}
\[
\vartheta=(M,K)\sim p(\vartheta), \qquad \vartheta\in\mathbb{R}_+^2
\]

\textbf{Known / fixed:}
network coupling $B_{ij}$ (IEEE-14), damping $D$, nominal injections $P_{m,i}$.

\vspace{0.4em}
\textbf{Planning-level control decision (fast frequency response):}
\[
u^{\text{ctrl}}_{\gamma,i}(t)=\gamma\,g_i\,\omega_i(t),
\qquad \sum_i g_i=1,\ g_i\ge0
\]

\textbf{Security-constrained requirement:}
\[
\gamma^*(\vartheta)=
\min_{\gamma}
\quad \text{s.t.}\quad
\max_t|\dot f(t)|\le r_{\max},\;
\min_t f(t)\ge f_{\min}.
\]

\textbf{Goal:}
design experiments that reduce uncertainty in the \emph{decision-relevant}
quantity $\gamma^*(M,K)$.

% =========================================================
\newpage
\section*{2. Finite-Horizon Sequential Active Probing Experiment}

\textbf{Episode structure:}
\[
\vartheta\sim p(\vartheta)\ \text{once and fixed},\qquad t=1,\dots,T.
\]

At step $t$, the experimenter chooses a probing design
\[
\xi_t=(b_t,A_t,T_p)\in\Xi=\mathcal{B}\times\mathcal{A}\times\{T_p\}.
\]

\textbf{Probe injection (active power, Hann window):}
\[
u^{\text{probe}}_{\xi,i}(\tau)=
\begin{cases}
A_t\,s(\tau;T_p), & i=b_t,\\
0, & i\neq b_t,
\end{cases}
\quad
s(\tau;T_p)=\frac12\!\left(1-\cos\frac{2\pi\tau}{T_p}\right).
\]

\textbf{Observation generation:}
simulate the swing equation under $(\vartheta,\xi_t)$, measure frequency,
compute observation $y_t$.

\textbf{Information available to the designer:}
\[
h_t=\{(\xi_1,y_1),\dots,(\xi_t,y_t)\},
\]
while $(M,K)$ remain unobserved.

% =========================================================

\section*{3. Active Probing Parameter Table (Fixed Design Choices)}

\begin{center}
\renewcommand{\arraystretch}{1.35}
\begin{tabular}{|c|c|c|p{4.6in}|}
\hline
\textbf{Parameter} & \textbf{Symbol} & \textbf{Value / Set} & \textbf{Justification (real references)} \\
\hline
Probe location &
$b_t$ &
$\mathcal{B}\subset\{1,\dots,14\}$ &
Buses with controllable active power (IBR-like) and good observability; standard in probing-based inertia and FR estimation.
\newline
Chakraborty et al., 2022; Peng et al., 2024. \\
\hline
Probe amplitude &
$A_t$ &
$\mathcal{A}=\{A_1,A_2,A_3\}$ &
Chosen so induced ROCOF exceeds PMU noise floor while remaining within security limits; validated in PHIL / probing studies.
\newline
Chakraborty et al., 2022; Yi et al., 2024. \\
\hline
Probe duration &
$T_p$ &
$2~\mathrm{s}$ &
Concentrates probing energy in inertial / primary-frequency band ($<0.5$ Hz) and avoids slower secondary dynamics.
\newline
Yi et al., 2024; Jiang et al., 2024. \\
\hline
Probe shape &
$s(t)$ &
Hann window &
Reduces spectral leakage and sharp transients; widely used in probing and system identification.
\newline
Yi et al., 2024. \\
\hline
Sampling rate &
$f_s$ &
$12~\mathrm{Hz}$ &
Consistent with PMU-class frequency measurements and ROCOF standards.
\newline
ENTSO-E, 2020; EPRI/NASPI, 2021. \\
\hline
Observation window &
$T_{\text{obs}}$ &
$[0,10]~\mathrm{s}$ &
Captures ROCOF peak and early frequency response, excluding secondary control.
\newline
ENTSO-E, 2020; Wang et al., 2022. \\
\hline
\end{tabular}
\end{center}


% =========================================================
\newpage
\section*{4. Observation Model and Feature Extraction (ROCOF-only)}

\textbf{PMU-like frequency measurement:}
\[
\Delta f_i(t)=\frac{\omega_i(t)}{2\pi}, \qquad t=n\Delta t,\ \Delta t=1/f_s.
\]

\textbf{General feature map (many possible):}
\[
\phi(\Delta f)=
[\mathrm{ROCOF}_{\max}, f_{\min}, t_{\text{nadir}}, t_{\text{settle}}, \dots].
\]

\textbf{ROCOF-only observation (used in this work):}
\[
\widehat{\dot f}(n)=
\frac{\Delta f((n+1)\Delta t)-\Delta f(n\Delta t)}{\Delta t},
\qquad
y_t=\mathrm{ROCOF}_{\max}=\max_{n\in\mathcal{W}_t}|\widehat{\dot f}(n)|.
\]

\textbf{Why ROCOF-only is valid:}
\begin{itemize}
\item ROCOF is directly linked to inertia and early frequency security margins.
\item It is widely used in inertia monitoring and probing-based estimation.
\item Low-dimensional observation reduces modeling error and stabilizes DAD training.
\end{itemize}

\textbf{References:}
ENTSO-E (2020); Wang et al. (2022); Chakraborty et al. (2022).

% =========================================================
\newpage
\section*{5. Likelihood Modeling and Why DAD (not iDAD)}

\textbf{Key clarification:}
the swing equation is deterministic; randomness enters via
measurement noise and unmodeled effects.
We define an explicit \emph{measurement-level} likelihood.

\textbf{Deterministic simulator mean:}
\[
\mu(\vartheta,\xi_t)=
\mathrm{ROCOF}_{\max}\big(\Delta f(\cdot;\vartheta,\xi_t)\big).
\]

\textbf{Explicit likelihood (offline training):}
\[
p(y_t\mid \vartheta,\xi_t)=
\mathcal{N}\!\left(\mu(\vartheta,\xi_t),\sigma^2\right).
\]

\textbf{Why DAD is selected:}
\begin{itemize}
\item DAD requires numerical evaluation of $\log p(y\mid\vartheta,\xi)$,
not an analytic physics likelihood.
\item An explicit, testable likelihood exists at the measurement level.
\item The iDAD paper explicitly recommends using DAD whenever an explicit
likelihood is available, as it serves as an upper bound.
\end{itemize}

\textbf{References:}
Foster et al. (2021); Ivanova et al. (2021).

% =========================================================
\newpage
\section*{6. DAD Training, MOCU, and End-to-End Workflow}

\textbf{Design policy (non-myopic):}
\[
\xi_t=\pi_\phi(h_{t-1}).
\]

\textbf{History update:}
\[
h_t=h_{t-1}\cup\{(\xi_t,y_t)\}.
\]

\textbf{Offline posterior (for evaluation):}
\[
p(\vartheta\mid h_T)\propto
p(\vartheta)\prod_{t=1}^T p(y_t\mid \vartheta,\xi_t).
\]

\textbf{Mean Objective Cost of Uncertainty (MOCU):}
\[
\mathrm{MOCU}(h_T)=
\mathbb{E}_{\vartheta\sim p(\vartheta\mid h_T)}
\big[\gamma^*(\mathcal{A}_T)-\gamma^*(\vartheta)\big].
\]

\textbf{Training objective (decision-aware DAD):}
\[
\phi^*=\arg\min_\phi\ \mathbb{E}[\mathrm{MOCU}(h_T)].
\]

\textbf{Workflow summary:}
\begin{itemize}
\item Offline: simulate episodes, evaluate likelihood, train $\pi_\phi$.
\item Online: apply $\pi_\phi$ sequentially using history only.
\item Baselines are myopic; DAD optimizes the \emph{entire probe sequence}.
\end{itemize}

\textbf{Reference:}
Foster et al. (2021).

